\documentclass[11pt,a4paper]{article}
\usepackage[T1]{fontenc}
\usepackage[utf8]{inputenc}
\usepackage{lmodern}
\usepackage[margin=24mm,headheight=14pt]{geometry}
\usepackage{amsmath,amssymb,amsthm,mathtools}
\usepackage{booktabs,longtable,array}
\usepackage{microtype}
\usepackage{xcolor}
\usepackage{xurl}
\usepackage{enumitem}
\usepackage{fancyhdr}
\usepackage[colorlinks=true,linkcolor=blue!45!black,citecolor=blue!45!black,
  urlcolor=blue!45!black,pdftitle={Saturated Pell certificates and quotient recovery for Diophantine consistency},
  pdfauthor={AI-assisted research contribution for the ProveIt project}]{hyperref}
\newcommand{\Z}{\mathbb Z}
\newcommand{\N}{\mathbb N}
\newcommand{\Zp}{\mathbb Z_{>0}}
\newcommand{\ConZFC}{\operatorname{Con}(\mathsf{ZFC})}
\newcommand{\RefZFC}{\operatorname{Ref}_{\mathsf{ZFC}}}
\newcommand{\file}[1]{\nolinkurl{#1}}
\newtheorem{theorem}{Theorem}[section]
\newtheorem{lemma}[theorem]{Lemma}
\newtheorem{proposition}[theorem]{Proposition}
\newtheorem{corollary}[theorem]{Corollary}
\theoremstyle{definition}
\newtheorem{definition}[theorem]{Definition}
\theoremstyle{remark}
\newtheorem{remark}[theorem]{Remark}
\numberwithin{equation}{section}
\setlist{itemsep=3pt,topsep=5pt}
\setlength{\parskip}{4pt}
\setlength{\parindent}{1.2em}
\setlength{\emergencystretch}{2em}
\pagestyle{fancy}
\fancyhf{}
\fancyhead[L]{\small Diophantine consistency: saturated certificates}
\fancyhead[R]{\small ProveIt research}
\fancyfoot[C]{\thepage}
\renewcommand{\headrulewidth}{0.3pt}
\allowdisplaybreaks[2]

\title{\vspace{-12mm}\textbf{Saturated Pell certificates and quotient recovery}\\[3mm]
  \large Eight explicit Diophantine equations for syntactic consistency of ZFC}
\author{\normalsize Research contribution prepared for Vladimir Reshetnikov's ProveIt project\\
  \small AI-assisted mathematical construction and exact independent audits}
\date{24 October 2026 (UTC)}

\begin{document}
\maketitle
\thispagestyle{plain}

\begin{abstract}
We give explicit, materialized integer polynomials with new variable--degree--cost
tradeoffs for the absence of a finite ZFC refutation. The main nine-positive
construction has exact total degree \(999\,518\); a shorter companion uses
435 binary arithmetic operations at degree \(1\,628\,768\). A strong-Pell
growth argument gives eleven positive unknowns at degree 213 in 112 operations.
For unrestricted integer unknowns, quotient recovery and a separated rational
norm give seventeen unknowns at degree 88 and sixteen at degree 184; a
degree-ten companion uses twenty-three integers.

The improvements are complete ordinary arithmetic reductions. They compose
with explicitly identified inherited source/compiler and logical-calibration
theorems; they do not assume consistency. Every fixed integer is supplied,
and the archive includes literal division-free circuits, exact algebra
certificates, ordinary proofs, and standard-library verifiers. The positive
majorant is certified by an identity with 36,309 positive coefficients.
We do not claim a global minimum, external historical priority, or a new
end-to-end proof-assistant derivation inside the encoded ZFC calculus.
\end{abstract}

\section{Outcome, target, and comparison}
\label{sec:outcome}

\subsection{The precise logical target}

Fix the repository's syntactic ZFC axiom predicate and proof calculus. Let
\(\RefZFC(p)\) mean that \(p\) codes a finite derivation of falsity, with
all axiom and inference checks included. We use
\[
 \ConZFC\ \equiv\ \neg\exists p\in\N\;\RefZFC(p).
\]
The desired conclusion for a fixed polynomial \(P\) and an explicitly
declared witness domain \(\mathcal D\) is
\begin{equation}
 \mathsf{ZFC}\ \vdash\
 \left[
   \left(\neg\exists\mathbf x\in\mathcal D\;P(\mathbf x)=0\right)
   \longleftrightarrow \ConZFC
 \right].
 \label{eq:target}
\end{equation}
This is a proof of an equivalence, not a proof of either side. In particular,
there is no appeal to the consistency or soundness of ZFC as a premise.
The relevant construction proves an existential equivalence with a finite
refutation and then negates it.

There is no single established numerical meaning of \emph{simplest}. A smaller
number of unknowns may require a higher degree; a lower degree may need a
larger arithmetic circuit; compact native constants can require a different
source packing. We therefore provide the measured alternatives in
Table~\ref{tab:results}, instead of asserting a nonexistent optimality theorem.
The distinction between positive and unrestricted integer witnesses is
substantive and is retained in every claim.

\begin{table}[htbp]
\centering
\small
\begin{tabular}{llrrr}
\toprule
Construction & Witness domain & Unknowns & Exact degree & Operations\\
\midrule
Complete-radical H1868 & \(\Zp\) & 9 & \(\mathbf{999\,518}\) & 1086\\
Simple H1868 & \(\Zp\) & 9 & \(1\,474\,688\) & 449\\
Simple compact H2082 & \(\Zp\) & 9 & \(1\,628\,768\) & \(\mathbf{435}\)\\
Strong-Pell weight & \(\Zp\) & 11 & \(\mathbf{213}\) & 112\\
Rational norm, product finish & \(\Z\) & 16 & \(\mathbf{184}\) & 135\\
Rational norm, squared finish & \(\Z\) & 16 & 192 & 134\\
Quotient recovery & \(\Z\) & 17 & \(\mathbf{88}\) & 114\\
Low-degree synthesis & \(\Z\) & 23 & \(\mathbf{10}\) & 131\\
\bottomrule
\end{tabular}
\caption{Complete new polynomials. All degrees count only the listed
unknowns, with the supplied coefficient tuple fixed.}
\label{tab:results}
\end{table}

For a readily inspectable equation over unrestricted integers,
\(P_{17}\) in Section~\ref{sec:integer17} is a useful primary presentation:
five squared residuals, seventeen unknowns, exact degree 88. If one fewer
integer unknown matters more than degree, Section~\ref{sec:integer16}
gives the sixteen-variable alternatives. If the number of positive unknowns
is the leading measure, Section~\ref{sec:nine-saturation} supplies the
nine-variable constructions. These choices do not displace shorter
higher-variable equations already present in the repository.

\subsection{Measured improvements over the inspected inputs}

The reviewed repository commit is
\[
 \texttt{c1a102936aca437e9999\allowbreak530a0111e6ddab56dc6d}.
\]
The review includes all five archives then present in \file{docs/incoming};
their exact byte hashes are recorded in \file{provenance/SOURCES.json}.
The comparisons below refer to these fixed inputs, not to a continually
changing branch.

\begin{table}[htbp]
\centering
\small
\begin{tabular}{p{43mm}p{49mm}p{48mm}}
\toprule
Comparison & Inspected input & New construction\\
\midrule
Nine positives, small circuit &
Degree \(1\,845\,940\), 440 operations \cite{repo} &
Degree \(1\,628\,768\), 435 operations\\
Nine positives, lower degree &
Degree \(1\,158\,228\), 1753 operations \cite{radical} &
Degree \(999\,518\), 1086 operations\\
Eleven positives, centered &
Degree 257, 114 operations \cite{joint} &
Degree 213, 112 operations\\
Integers, fixed degree 88 &
19 unknowns, 109 operations \cite{joint} &
17 unknowns, 114 operations\\
Integers, fixed degree ten &
24 unknowns \cite{budget} &
23 unknowns, 131 operations\\
\bottomrule
\end{tabular}
\caption{Specific comparisons. A variable reduction can increase the
operation count; it is not described as improvement in every measure.}
\label{tab:comparisons}
\end{table}

The incoming standalone report also gives a nine-positive degree
\(1\,158\,240\) equation in 811 operations \cite{reportincoming}.
The new degree-\(999\,518\) form lowers its reported degree but uses more
operations. That standalone report contains no executable companion in its
archive, so its operation figure is treated as a reported comparator.
The incoming eleven-positive degree-277 / 111-operation equation remains
one operation shorter than our degree-213 form \cite{joint}.
The active size plan also retains a fifteen-positive degree-157 equation
in only 66 operations \cite{repo}; we make no claim to beat that total
operation count.

\subsection{Counting rules and fixed coefficients}

The operation measure counts each live binary addition, subtraction, and
multiplication in the delivered straight-line circuit once. A square is
one multiplication. Multiplication by a fixed scalar is also counted.
Intermediate registers can be reused. Fixed integer leaves are free in
this arithmetic-circuit measure, so their storage cost is stated separately.
No division, variable exponent, inequality test, logarithm, bit operation,
or proof predicate occurs in a delivered polynomial evaluator.

The following are the native coefficient costs, not the coefficients of a
fully expanded enormous polynomial:
\begin{center}
\small
\begin{tabular}{lrrr}
\toprule
Coherent tuple & Ports & Raw bytes & Largest port, bits\\
\midrule
Compact54-sharp (nine, H2082) & 4 & 215492 & 434054\\
Original31 (nine, H1868) & 4 & 533099 & 1070004\\
Compact44 (eleven positives) & 4 & 295388 & 592884\\
Source34 plus determinant scales & 6 & 615883 & 1065043\\
\bottomrule
\end{tabular}
\end{center}
The full raw and compressed costs, individual bit lengths, and raw hashes
are in \file{COEFFICIENT_COSTS.json}. The nine-variable circuits also
materialize derived kernel coefficients and power-of-two scale leaves;
their bindings and actual bytes are supplied separately. These are
polynomials with specified finite integer constants, not a universal
polynomial with an unspecified program-index parameter. The coherent
tuples must not be mixed merely because two radix exponents happen to agree.

For witnesses in \(\N=\{0,1,\ldots\}\), any positive-domain construction
can be translated by \(x_j=u_j+1\). This preserves total degree and
existence equivalence. The reported operation counts are for the original
positive circuits; the translated circuit is a different counted artifact.
Likewise, the integer constructions are not obtained by silently allowing
negative inputs in a positive-domain theorem.

\section{Proof architecture and inherited dependencies}
\label{sec:status}

\subsection{What the ordinary proof composes}

Each component establishes, in its stated domain,
\[
 \exists\mathbf x\;P_{\rm new}(\mathbf x)=0
 \quad\longleftrightarrow\quad
 \exists\mathbf y\;\mathcal S_{\rm authenticated}(\mathbf y)=0.
\]
The authenticated arithmetic sources have inherited ordinary theorems
relating their zeros to the fixed native computation and then to
\(\exists p\,\RefZFC(p)\). Composition gives
\eqref{eq:target}. The archive retains the relevant source proofs,
the concrete coefficient adapters, and their provenance.
It does not re-prove the entire earlier machine/tableau/calculus chain.

In the nine-positive branch the inherited interfaces are the native
positive source, the exact sparse carry packing, the numerical ratio
theorem, and the old ProductPell certificate. We prove the new saturated
certificate and both complete finalizers. The eleven-positive branch
inherits a weighted native Pell source and its canonical weight lifting;
we prove the smaller weight works before the required source positivity
has been established. The integer branch inherits signed orientation,
strong-rank recovery, and canonical determinant completeness; we prove
the quotient elimination and subsequent norm merger without assuming
either omitted equation in their own justifications.

These are ordinary unbounded proofs. The hypothetical proof or computation
can be arbitrarily long. The arguments use finite polynomial identities,
integer divisibility, induction on recurrence indices, and standard
Pell and three-square results available in ZFC. They never infer a
universal result from a bounded search.

\subsection{Three different verification artifacts}

It is necessary to distinguish:
\begin{enumerate}
 \item An ordinary mathematical proof of the new reductions, composed
 with the named ordinary source/calibration lemmas.
 \item Exact finite verification of the delivered identities, constants,
 circuits, source restorations, and degree certificates.
 \item A proof-assistant theorem for the entire new equation, or an
 actual derivation inside the repository's encoded object-language ZFC.
\end{enumerate}
This contribution supplies the first two. It does not supply the third.
The exact K215 machine already has external Lean and Rocq correctness
theorems \cite{repo}; the recent native arithmetic/compiler and internal-ZFC
formalization have distinct remaining obligations. In particular, an
external theorem about the standard arithmetic predicate is not
automatically an exported object-language proof covering every internal
natural of a nonstandard model.

The retained \file{provenance/OBJECT_ZFC_AT_PIN.md} states that boundary
precisely. We do not promote the finite Python audits, or a general claim
of formalizability, into a newly checked internal biconditional. The
ordinary theorem in this report is a composition with identified
inherited mathematical lemmas; it is not a fresh end-to-end kernel
certificate. The repository's current size-reduction plan explicitly
pauses extension of Lean formalization \cite{repo}.

\subsection{Scope of the new ideas}

The incoming round-11 work already used norm saturation for a different
ten-integer Sun/Lucas construction \cite{round11}. Incoming papers also
used whole-radical approximation, root centering, determinant guards,
joint positivity, and adjacent-coefficient rational Pell descent
\cite{radical,budget,joint}. Those ideas are credited.

The new content is the positive ProductPell saturation theorem, its
exact full-radical correction and measured complete equations; the
strong-Pell growth bound replacing a larger guard by \(i^3+cT\);
the explicit removal of a quotient witness after a noncircular bootstrap;
and the separated weighted rational norm that removes one further
integer coordinate. The lower-degree twenty-three-integer equation
is a new composition of already available local transformations.
These are precise improvements relative to the reviewed sources.
No claim of historical priority beyond that comparison is made.

\section{Positive Pell saturation and a degree below one million}
\label{sec:nine-saturation}

This section gives three explicit polynomials in the same nine strictly
positive unknowns
\[
 b,g,h,s,w,C,i,j,n\in\Zp.
\]
Their other registers are polynomial abbreviations, not additional
existential variables. The finite source tables, fixed integer payloads,
and division-free arithmetic circuits accompany the report.
\begin{center}
\begin{tabular}{lrr}
\hline
Polynomial & Exact degree & Binary operations\\
\hline
Compact H2082, simple majorant & \(1\,628\,768\) & \(435\)\\
H1868, simple majorant & \(1\,474\,688\) & \(449\)\\
H1868, complete-radical majorant & \(\mathbf{999\,518}\) & \(1086\)\\
\hline
\end{tabular}
\end{center}
An operation is one addition, subtraction, or multiplication; squaring and
scalar multiplication each cost one, shared registers are reused, and fixed
integer leaves cost zero. These are different points in a degree--cost
comparison. The first improves the inspected 440-operation H2082 equation;
the last improves the inspected nine-positive degree frontier, including the
incoming radical-guard construction \cite{repo,radical,reportincoming}.
The variable count is unchanged.

The incoming round-11 package already applies \emph{norm saturation} to a
different ten-integer Sun/Lucas certificate \cite{round11}. We credit that
idea explicitly. The new steps here are its transfer to the positive
ProductPell certificate, a majorant of the complete radical \(DFI\), and their
composition with the H1868 source. The whole-radical approximation viewpoint
also builds on the incoming majorant work \cite{radical,reportincoming}.
No historical priority or global minimum is claimed.

\subsection{Removing one square factor from the positive certificate}

For \(A\ge2\), let
\[
 \psi_A(0)=0,\qquad \psi_A(1)=1,\qquad
 \psi_A(t+2)=2A\psi_A(t+1)-\psi_A(t),
\]
and let \(\chi_A\) have initial values \(1,A\) and the same recurrence.
Thus
\[
 \chi_A(t)^2-(A^2-1)\psi_A(t)^2=1.
\]
We use the standard positive Pell parametrization and the divisibility and
congruence lemmas used in Jones's ProductPell proof \cite{jones,productpell}.

\begin{lemma}[Saturated positive ProductPell]
\label{lem:saturated-positive}
Let \(A\ge2\), \(B\ge3\) be odd, \(C\ge B\), and \(J\ge1\).
For \(i,j\in\Zp\), define
\begin{align}
 D&=(A^2-1)C^2+1,& r&=iJC^2,\notag\\
 F&=(A^2-1)Dr^2+1,& L_m&=CD+1-A,\notag\\
 G&=A+FL_m,& H&=C+F\bigl(B+(j-1)C\bigr),\notag\\
 I&=(G^2-1)H^2+1. \label{eq:saturated-pell}
\end{align}
Then
\[
 C=\psi_A(B)
 \quad\Longleftrightarrow\quad
 \exists i,j\in\Zp\;\;DFI\text{ is an integer square}.
\]
\end{lemma}

\begin{proof}
All displayed radicands are positive. The defining equations give
\[
\begin{array}{lll}
 F\equiv1\pmod D,&G\equiv1\pmod D,&I\equiv1\pmod D,\\
 F\equiv1\pmod C,&G\equiv1\pmod C,&H\equiv B\pmod C,\\
 G\equiv A\pmod F,&H\equiv C\pmod F,&I\equiv D\pmod F.
\end{array}
\]
Consequently \(D,F,I\) are pairwise coprime. A square product therefore
makes each factor a square. Write \(D=d^2\), \(d>0\), and introduce only
inside the proof
\[
 E'=iJdC^2.
\]
Then \(F=(A^2-1)E'^2+1\) and \(C^2\mid E'\). Pell parametrization supplies
\[
 C=\psi_A(t),\qquad E'=\psi_A(u),\qquad H=\psi_G(v).
\]
Here \(2\le t<C\), and the square-divisibility lemma
\[
 \psi_A(t)^2\mid\psi_A(u)\ \Longrightarrow\ \psi_A(t)\mid u
\]
gives \(C\mid u\). Since \(F=\chi_A(u)^2\), recurrence congruence and
\(H\equiv C\pmod F\) imply
\[
 \psi_A(v)\equiv\psi_A(t)\pmod{\chi_A(u)}.
\]
The Pell step-down congruence gives \(v\equiv\pm t\pmod{2u}\), and hence
\(v\equiv\pm t\pmod C\). But \(G\equiv1\pmod C\) gives
\(\psi_G(v)\equiv v\pmod C\), so \(B\equiv\pm t\pmod C\).
The bounds leave \(B=t\) or \(B+t=C\). The latter is impossible:
\(\psi_A(t)\equiv t\pmod2\) makes \(C-t\) even, whereas \(B\) is odd.

For completeness suppose \(C=\psi_A(B)\), and put \(d=\chi_A(B)\).
Periodicity of the invertible Pell recurrence modulo \(JDC^2\) supplies
a positive index \(u\) and \(i_{\rm old}>0\) with
\[
 \psi_A(u)=i_{\rm old}JDC^2.
\]
Set \(i=d\,i_{\rm old}\). Then the new \(F\) equals
\[
 (A^2-1)(i_{\rm old}JDC^2)^2+1=\chi_A(u)^2.
\]
Put \(H_0=\psi_G(B)\). Recurrence congruences give
\(H_0\equiv C\pmod F\) and \(H_0\equiv B\pmod C\).
Because \(F\equiv1\pmod C\), the integer \(q=(H_0-C)/F\) satisfies
\(q\equiv B\pmod C\). Moreover \(L_m\ge C\), \(G\ge FC\), and
\(\psi_G(B)\ge BG\) for \(B\ge2\); thus
\[
 H_0\ge BFC\ge FB+C,\qquad q\ge B.
\]
Therefore \(j=1+(q-B)/C\) is positive and realizes the displayed \(H\).
Finally \(I=\chi_G(B)^2\), proving the square-product condition.
\end{proof}

The earlier positive construction used
\[
 F_{\rm old}=(A^2-1)(i_{\rm old}JDC^2)^2+1.
\]
Its proof used \(D\mid E_{\rm old}\) only to obtain \(F\equiv1\pmod D\).
The new polynomial provides that congruence directly, while retaining
\(C^2\mid E'\), the divisibility actually needed by the index argument.
The completeness map \(i=d\,i_{\rm old}\) is an embedding, not a claimed
bijection of all witness tuples.

\subsection{The H1868 source as a positive-domain prefix}

The source data are pinned to repository snapshot
\texttt{c1a102936aca} (the full hash is in the provenance files) and the incoming
round-11 H1868 certificate \cite{repo,round11}. They retain the complete
original31 coefficient tuple. Its four fixed integers are materialized in
the archive; they are neither free parameters nor unspecified program
indices. The compact H2082 companion instead retains its own coherent
compact tuple from the source-54 lineage \cite{source54}.

For H1868 the 31 coordinate weights, in the supplied source order, are
\[
\begin{split}
&(1,7,61,36,16,85,20,67,43,11,34,13,96,80,29,50,\\
&\hspace{12mm}2,109,110,112,116,91,53,55,27,101,92,9,107,37,60).
\end{split}
\]
The 22 target positions, in source-row order, are
\[
\begin{split}
&(1141,28,1437,1069,1836,1354,1814,614,267,1408,1403,\\
&\hspace{12mm}889,337,1587,413,673,1139,772,1131,1018,1055,85).
\end{split}
\]
The literal files
\texttt{\detokenize{nine/data/source.json}},
\texttt{\detokenize{nine/data/code.json}}, and
\texttt{\detokenize{nine/data/kernels.json}}
specify the source residuals \(\rho_j\) and the three integer polynomials
\(D_0,D_1,D_2\). They satisfy, as exact polynomial identities,
\begin{equation}
 [T^{q_j}]\left(D_2(T)U_{\mathbf p}(T)^2+
 D_1(T)U_{\mathbf p}(T)+D_0(T)\right)=2\rho_j(\mathbf p),
 \qquad U_{\mathbf p}(T)=\sum_\ell p_\ell T^{w_\ell}.
 \label{eq:h1868-identities}
\end{equation}
Equal kernel coefficients at a shared position are retained once.
The audit checks all absent monomials as well as the intended ones:
22 identities, 20,466 convolution interactions, and 81 stored positions.
The kernel degrees are \((1408,1727,1636)\), with leading coefficients
\((-6,2,1)\); therefore the convolution height is \(1868\), the maximum
weight is \(116\), and the maximum target is \(1836\).

Put \(\kappa=2134939\). Define
\begin{align}
 B_c&=2^\kappa b^2,&
 \ell&=\sum_\ell B_c^{w_\ell},&
 T_c&=\sum_j B_c^{q_j},\notag\\
 U_s&=\ell+g,&
 G_s&=(D_2(B_c)U_s+D_1(B_c))U_s+D_0(B_c),\notag\\
 N_0&=B_c^{117},&
 M_c&=N_0-1-(b-1)\ell,\notag\\
 S_c&=g+N_0(2G_s+T_c),&
 T_m&=M_c+N_0(B_c-2)T_c,\notag\\
 L_c&=2bB_c^{1986},&
 R&=L_c(S_c+T_m+1)-S_c-1,\notag\\
 N&=2^{16012032}B_c^{1055},&
 t_c&=bB_c^{116}. \label{eq:h1868-prefix}
\end{align}
The factor \(2\) in \(L_c\), and the absence of an additional \(b\) in
\(N\), are essential.

Here every kernel coefficient has magnitude below \(2^{2134927}\).
The quadratic layer dominates the linear and constant layers by 25 and
460 powers of \(B_c\), respectively. The inherited geometric-series and
sparse-carry bounds consequently give \(G_s>0\), \(R,N\ge8\), and
\(b\le R,N\) on all positive inputs, before the base is known to be binary.
On the recovered range \(g<4t_c\), they give
\[
 S_c+T_m<B_c^{1986}<L_c.
 \label{eq:h1868-range}
\]
The coefficient bound and the unchanged 32-position target headroom are
exactly the hypotheses used by the inherited sparse-carry proof
\cite{round11,source71}.

For clarity, the binary threshold can be recalculated directly.
If \(b=2^k\), \(B_c=2^v\), and \(L_c=2^u\), then
\[
 v=\kappa+2k,\quad u=1+k+1986v,\quad
 \operatorname{pop}(T_m)=139v-31k-22.
\]
The two-block identity for \(R\) gives
\[
 \operatorname{pop}(R)=u+\operatorname{pop}(T_m)-\tau,\qquad
 \tau=\operatorname{pop}(S_c)+\operatorname{pop}(T_m)
       -\operatorname{pop}(S_c+T_m)\ge0.
\]
Moreover
\[
 u+\operatorname{pop}(T_m)=2125\kappa+4210k-21=2\log_2N.
\]
Legendre's identity therefore gives
\begin{equation}
 N^2\mid \binom{2R}{R}
 \quad\Longleftrightarrow\quad S_c\mathbin{\mathrm{AND}}T_m=0.
 \label{eq:h1868-calibration}
\end{equation}
The sparse-carry theorem decodes the right side into a positive common
source zero. Conversely a fixed source zero packs for all sufficiently
large binary \(b\), with \(0<g<t_c\). Positivity of \(g\) follows because
the all-one source tuple is not a zero: its fifth residual equals \(1\).
The endpoint \(b=1\) cannot create a spurious solution, since its low mask
would force \(g=0\).

\subsection{The common ratio registers}

From either authenticated prefix define
\begin{align}
 U&=N^2w,&Y&=N^2s,&M&=RY,\notag\\
 P&=2M^2U,&K&=R+1+h(P-1),&A&=M(U+1),\notag\\
 B_1&=2R+1,&W_1&=bw,&J&=4A-5,\notag\\
 Q&=3W_1C-2(W_1^2-1),&
 a_0&=(P^2-1)K^2+1,\notag\\
 m&=K^2(4t_c-g)-16t_c(C-KY)^2,&
 \Omega&=Jm,\notag\\
 E_0&=J(n+W_1^2)+Q,&V_0&=PK. \label{eq:nine-ratio}
\end{align}
Use \eqref{eq:saturated-pell} with \(B=B_1\), and abbreviate
\[
 Z=B_1+(j-1)C,\qquad a_{1,\rm raw}=DFI.
\]
The inherited numerical ratio theorem converts the source conditions into
\[
 a_0\text{ square},\qquad C=\psi_A(B_1),\qquad
 J\mid Q,\qquad m>0,
 \label{eq:nine-ratio-contract}
\]
with the weak soundness bound \(g<4t_c\) and the stronger completeness bound
\(g<t_c\) \cite{jones,source71,productpell}.
When that theorem is phrased using the old ProductPell witnesses, the
equivalence is interpreted existentially. After recovering \(C=\psi_A(B_1)\)
we may reconstruct the old witnesses. We do not divide an arbitrary new
\(i\) by \(\sqrt D\).

\subsection{Two exact majorants}

The simple choice is
\begin{equation}
 T_{\rm simple}=ADr,\qquad V_1=T_{\rm simple}GH,\qquad
 a_1=a_{1,\rm raw}.
 \label{eq:simple-majorant}
\end{equation}
Its unconditional certificate is particularly short:
\begin{align}
 T_{\rm simple}^2-DF&=D(Dr^2-1)>0,\notag\\
 V_1^2-a_{1,\rm raw}
 &=D(Dr^2-1)G^2H^2+DF(H^2-1)>0.
 \label{eq:simple-majorant-certificate}
\end{align}
This supplies the 435-operation compact-H2082 polynomial and the
449-operation H1868 polynomial in the table.

For the lower-degree construction we approximate the \emph{complete}
radical \(\sqrt{DFI}\). Here is an exact finite definition of the supplied
333-term polynomial, so its construction does not depend on a numerical
approximation. Regard
\[
 \mathcal R(A)=D(A)F(A)I(A)=\sum_{k=0}^{36}q_k(C,i,Z)A^k
\]
as a polynomial in \(A\), with \(J=4A-5\).
Its leading coefficient is \(2^{20}C^{38}i^{10}Z^2\). In the rational
function field \(\mathbb Q(C,i,Z)\), set
\begin{equation}
 t_{18}=1024C^{19}i^5Z,\qquad
 t_{18-r}=
 \frac{q_{36-r}-\sum_{s=1}^{r-1}t_{18-s}t_{18-r+s}}{2t_{18}}
 \quad(1\le r\le18).
 \label{eq:radical-recurrence}
\end{equation}
Put
\begin{align}
 \sigma&=1024C^5i,\notag\\
 \mathcal P_0&=\sigma\sum_{k=0}^{18}t_k A^k,\notag\\
 \mathcal P&=\mathcal P_0+297874C^{24}i^6Z.
 \label{eq:complete-majorant}
\end{align}
The finite coefficient calculation proves
\(\mathcal P_0,\mathcal P\in\mathbb Z[A,C,i,Z]\); the authoritative literal
integer table is
\texttt{\detokenize{nine/data/saturated_corrected_P.json}}.
The divisions in \eqref{eq:radical-recurrence} describe a finite way to
derive that fixed table. The delivered evaluator uses only its integer
coefficients, addition, subtraction, and multiplication.

\begin{proposition}[Complete-radical certificate]
\label{prop:complete-majorant}
For \(A\ge3\), \(C,i,Z\ge1\),
\[
 \mathcal P>\sigma\sqrt{DFI}>0.
\]
\end{proposition}
\begin{proof}
The independent exact checker reconstructs \(DFI\) from its defining
polynomials and verifies
\[
 \left(\mathcal P^2-\sigma^2DFI\right)
 (x+3,y+1,u+1,v+1)=\mathcal C(x,y,u,v),
\]
where the supplied coefficient table for \(\mathcal C\) has 36,309
strictly positive nonzero coefficients and constant term
\[
 8341416066461366340.
 \label{eq:defect-positive-constant}
\]
It also verifies that \(\mathcal P(x+3,y+1,i,Z)\) has positive
coefficients and a positive constant term. Thus the defect and the
majorant itself are strictly positive on the required domain. These exact
polynomial identities prove the inequality.
\end{proof}

For this variant use
\begin{equation}
 a_1=\sigma^2DFI,\qquad V_1=\mathcal P.
 \label{eq:complete-variant}
\end{equation}
The finite recurrence, literal coefficient tables, and positive-defect
certificate serve different purposes: the recurrence explains the
construction; the table defines the exact polynomial; the certificate
proves its uniform inequality. The Laurent approximation itself is not
used as an asymptotic assumption.

\subsection{The single equation and its proof}

For either choice \eqref{eq:simple-majorant} or
\eqref{eq:complete-variant}, put
\begin{equation}
 \Theta=E_0-\Omega(V_0+V_1),\qquad
 \boxed{\;
 \mathcal F_9=
 \left(\Theta^2-\Omega^2(a_0+a_1)\right)^2
       -4\Omega^4a_0a_1=0.\;}
 \label{eq:saturated-nine-equation}
\end{equation}

\begin{theorem}
\label{thm:saturated-nine}
The polynomial \(\mathcal F_9\) has a positive integer zero if and only if
its authenticated positive source has a zero. With the inherited
source-calculus calibration,
\[
 \mathrm{ZFC}\vdash
 \left(
   \neg\exists(b,g,h,s,w,C,i,j,n)\in\Zp^9\;
       \mathcal F_9=0
 \right)
 \longleftrightarrow\ConZFC.
\]
\end{theorem}

\begin{proof}
The unconditional source bounds give \(A\ge33280\), \(A>W_1\), and
\(J>3A>0\). Also
\[
 A^2-P=M^2(U^2+1)\ge1,\qquad
 D\ge A^2,\qquad D>C^2.
\]
Both radicands are positive. We have \(V_0^2-a_0=K^2-1>0\);
the chosen majorant gives \(V_1>\sqrt{a_1}\). Furthermore
\[
 E_0=Jn+(J-2)W_1^2+3W_1C+2>0.
\]
If \(m=0\), then \(\mathcal F_9=E_0^4>0\). Otherwise factoring the
two-root norm over the reals gives signs \(\epsilon_0,\epsilon_1\) with
\[
 \frac{E_0}{J}
 =m\left(V_0+V_1+\epsilon_0\sqrt{a_0}
                    +\epsilon_1\sqrt{a_1}\right).
\]
The parenthesis is positive for every sign pair. Thus \(m>0\)
\emph{before} any Pell-index or square-root-integrality argument.

The margin now gives \(|C-KY|<K/2\), and therefore
\[
 C>2K>2R+1+W_1>B_1.
\]
Since \(D>PC^2\), \(F>D\), and \(C>K\),
\[
 a_{1,\rm raw}>D^2>P^2K^2\ge a_0.
\]
This also separates the radicands in the scaled variant. Hence
\(\Theta/\Omega\), a signed sum of their roots, is rational and nonzero.
For signed roots \(u_0,u_1\),
\[
 2(u_0+u_1)u_0=(u_0+u_1)^2+a_0-a_1
\]
makes both roots rational, then integral. In the complete variant,
\(\sqrt{a_1}/\sigma\) is a rational square root of the integer
\(DFI\), so it too is integral. The same norm-factor identity yields
\(J\mid E_0\), and hence \(J\mid Q\).
Lemma~\ref{lem:saturated-positive} gives \(C=\psi_A(B_1)\).
The ratio theorem and the authenticated carry test now recover a source
zero. If necessary, reconstruct the old ProductPell witnesses before
invoking the original ratio theorem; the other ratio registers remain
fixed.

Conversely a source zero supplies the ratio conditions, with \(m>0\).
Lemma~\ref{lem:saturated-positive} supplies positive \(i,j\) making \(DFI\)
a square. Let \(z_0,z_1\) be the positive integral roots of the chosen
\(a_0,a_1\), and define
\[
 n=m(V_0+V_1+z_0+z_1)-W_1^2-Q/J.
\]
This is an integer. It is positive because
\[
 W_1^2+Q/J
 =(1-2/J)W_1^2+(3W_1C+2)/J<2D,
\]
whereas \(V_1>4D\). For the simple majorant the latter follows directly
from \(V_1=ADrGH\). For the complete majorant,
\(F>3D\), \(G\ge3\), and \(H\ge2\) give \(DFI>99D^2\), and
Proposition~\ref{prop:complete-majorant} applies. This choice of \(n\)
makes one norm factor vanish. The existence equivalence follows, and
negating it gives the displayed consistency equivalence.
\end{proof}

\subsection{Why the degree is exactly 999,518}

The complete-radical construction needs a cancellation-aware calculation.
For H1868 the positive leading terms of the source are
\[
 R\sim4bB_c^{3971},\qquad
 N=2^{16012032}B_c^{1055}.
\]
Thus, writing \(r_R=\deg R\), \(\nu=\deg N\), and \(a=\deg A\),
\[
 r_R=7943,\qquad \nu=2110,\qquad
 a=r_R+4\nu+2=16385.
\]
For weights \((a,1,1,r_R)\) on \((A,C,i,Z)\), the unique leading
monomials of the complete majorant and its defect are
\[
\begin{array}{ll}
 \mathcal P:&1048576A^{18}C^{24}i^6Z,\\
 \mathcal P^2-\sigma^2DFI:
   &624687054848A^{18}C^{48}i^{12}Z^2 .
\end{array}
\]
Consequently
\[
 \deg V_1=18a+r_R+30=302903,\qquad
 \deg\delta_1=18a+2r_R+60=310876,
\]
where \(\delta_1=V_1^2-a_1\). Also, for \(\delta_0=V_0^2-a_0\),
\[
\begin{aligned}
 \deg\Omega&=5a+2W+3,&\deg E_0&=a+4,\\
 \deg V_0&=3a+r_R+1,&\deg\delta_0&=3a+r_R+2.
\end{aligned}
\]
Put
\[
 \mathcal B=E_0^2-2E_0\Omega(V_0+V_1)
                         +\Omega^2(\delta_0+\delta_1).
\]
Its unique leading contribution is \(\Omega^2\delta_1\). The exact identity
\[
 \mathcal F_9=\mathcal B^2+4\Omega^2V_0V_1\mathcal B
            +4\Omega^4(\delta_0V_1^2+a_0\delta_1)
\]
has four displayed contribution degrees
\[
 950392,\qquad \boxed{999518},\qquad 991546,\qquad 753714.
\]
The maximum is unique and its leading coefficient is positive. Therefore
\[
 \deg\mathcal F_9=59a+4r_R+8W+103=999518
\]
is attained, not merely an upper bound from a circuit.

\paragraph{Verification and proof boundary.}
The independent standard-library audit reconstructs the source identities,
the positive population threshold, the 3035-term local radicand, the
333-term majorant, the complete positive-defect certificate, and the exact
degree comparison without importing the producer. A separate literal
circuit checker verifies the exported equations and operation counts.
The three complete Pell fixtures and 5472 local arithmetic checks support
transcription checks; they are not experiments concerning ZFC's
consistency. The consistency theorem uses the identified ordinary
source/compiler calibration \cite{repo,source71}. This work supplies no new
end-to-end Lean/Rocq certificate or exported derivation in an internal ZFC
calculus.

\section{An eleven-positive equation of degree 213}
\label{sec:strong-pell-budget}

This companion lowers the incoming centered eleven-positive equation from
degree $257$ and $114$ operations to degree $213$ and $112$ operations.
The incoming degree-$277$ equation still has one fewer operation
\cite{joint}.  The new mechanism is to use growth forced by the strong
Pell equation to bound the transport guard, then center the norm at
\emph{both} the strong and auxiliary roots.  All statements in this
section concern ordinary proofs formalizable in ZFC.  The compact44
native compiler and the normalized source calibration are inherited;
this is not a new end-to-end proof-assistant certificate.

\subsection{The complete polynomial}

The unknowns are
\[
 (j,F,\alpha,i,T,X,Y,n,\eta,\zeta,\gamma)\in\Zp^{11}.
\]
Use the entire compact44 coefficient tuple
$(\beta,K_0,M_C,M_F)$, with $\beta=2^{590069}-1$ and the already shifted
$M_F$.  The companion includes their actual integer payloads and binding;
these ports are not interchangeable with a different native tuple.
Every following name abbreviates a polynomial and adds no unknown:
\begin{align*}
 J&=3j+1,&q&=\beta J+1,&C&=q-F-\alpha,&E&=XY,\\
 k&=\eta+\zeta,&c&=kY+\eta,&a&=E+Y,&H&=4a+3,\\
 \Delta&=a^2+H,&\ell&=EkY,&I&=ic^2,&D_f&=1+\Delta I^2,\\
 Q&=\Delta^2I^2,&\delta&=Q-1,&M&=q^3(q-1).
\end{align*}
Further put
\begin{align*}
 R&=((q-1)(q-F)+\alpha)(q^2-1)+(M_C+qM_F)J,\\
 B_t&=(qK_0+X)C-qF,&D_\tau&=\ell(\ell+k)+1,\\
 V&=D_f(cT-R)-c,&\bar f&=(a+2)I,\\
 h_*&=i^3+cT,&W&=1+(q-1)h_*,\\
 \bar Z&=M(\ell+k+n+cT+V+\bar f)+(q-1)Y-q^3WB_t.
\end{align*}
The unchanged main carrier is
\begin{align*}
 v&=k-R-1,&S&=4qY,&U&=qS,\\
 D_*&=U(X+ac)+H(U\gamma+SX-v),\\
 A_1&=D_*^2-U^2(1+\Delta c^2).
\end{align*}
The octic norm is evaluated by
\begin{align*}
 A_r&=M^2D_\tau,&B_r&=M^2D_f,&C_r&=M^2(QV^2-1),\\
 T_N&=\delta^2A_rB_r,&R_N&=\delta C_r,\\
 J_N&=\delta(\bar Z^2-A_r-B_r)-C_r,\\
 N_c&=J_N^2-4\{T_N+R_N(A_r+B_r)\},\\
 N_8&=N_c^2-64T_NR_N\bar Z^2.
\end{align*}
The single equation is
\begin{equation}
 \boxed{\quad P_{11}=N_8+A_1+N_8A_1=0.\quad}
 \label{eq:pell-budget11}
\end{equation}
It uses only fixed integer coefficients, fixed exponents, addition,
subtraction, and multiplication.  Since $4\mid U$, both $A_1$ and $N_8$
are squares modulo $4$.  Thus $(1+N_8)(1+A_1)=1$ forces both factors to
be $1$: the alternative $(-1,-1)$ would require residues $-2$.
Consequently \eqref{eq:pell-budget11} is exactly the conjunction
$N_8=A_1=0$ on its stated domain.

\subsection{Source boundary and recovery before positivity}

The inherited weighted source theorem \cite{source70,joint} concerns
positive coordinates satisfying the main residual, the strong norm,
the auxiliary norm, and the first carrier
\begin{align*}
 f^2-\Delta(ic^2)^2&=1,\\
 QV^2-(Q-1)(V+u)^2&=1,\\
 Z^2&=M^2D_\tau,\\
 Z&=M(\ell+k+r)+(q-1)Y-q^3WB_t.
\end{align*}
Its soundness holds for every positive integer $h_*$.  For canonical
witnesses, after the strong and auxiliary pairs have been chosen, a
chosen weight can be lifted by
\[
 r=Wt_0+\tau-\ell-k-s>qW-2k,
 \qquad t_0\ge q,\quad 0<s<k,\quad \ell<\tau<\ell+k.
\]
The source has a positive zero precisely when the fixed ZFC calculus has
a finite refutation.  The authentic mask bounds are
\[
 K_0\ge1,\quad 0<M_C<\beta,\quad0<M_F<2\beta,
 \qquad 0<(M_C+qM_F)J<2q^2.
\]
Once $C>0$, these imply $3q<R<q^4-2$.

Here is the part of the norm recovery that must precede $C>0$.
For every positive input, $I\ge c^2$, $\Delta\ge15$, and
\[
 1<D_\tau<D_f<Q,\qquad |V|>1.
\]
The rational radicand
$D_y=(QV^2-1)/(Q-1)>V^2$ is distinct from $D_\tau,D_f$.
Moreover, with $t=\sqrt{D_\tau}$, $f=\sqrt{D_f}$,
$y=\sqrt{D_y}$, no signed sum of $t,f,y$ is zero.
If $cT-R=0$, then $f>t+y$; otherwise $y>f+t$.
These inequalities depend only on $I\ge c^2$, not on $C$ or a bound on
$R$; their elementary proof is retained in the companion source audit.

The universal three-root norm therefore gives, at $N_8=0$,
\begin{equation}
 \bar Z/M=\tau_*+f_*+y_*,\qquad
 \tau_*^2=D_\tau,\quad f_*^2=D_f,\quad y_*^2=D_y.
 \label{eq:budget-radical-recovery}
\end{equation}
All three roots are rational.  One can prove this without assuming
independent square classes: for signed roots $x_1,x_2,x_3$ with rational
nonzero sum $t_0$, put
\[
 u_0=(t_0^2-x_1^2-x_2^2-x_3^2)/2,
 \qquad
 v_0=\frac{u_0^2-x_1^2x_2^2-x_1^2x_3^2-x_2^2x_3^2}{2t_0}.
\]
Then $u_0$ and $v_0$ are rational and
$x_1=x_1^2(x_1^2+u_0)/(x_1^2t_0+v_0)$.
The denominator is $x_1(x_1+x_2)(x_1+x_3)\ne0$, by distinctness of
the radicands; use the cyclic formulas as well.
Thus $\tau_*,f_*$ are integers.  The adjacent-coefficient Pell descent
from $QV^2-(Q-1)y_*^2=1$ makes $y_*$ an integer too
\cite{joint}.  In particular $M\mid\bar Z$.

Reduction modulo $q^3$ and $q-1$ now gives $q^3\mid Y$ and
$q-1\mid B_t$, independently of the sign of $C$.
The main residual makes $g=|D_*/U|$ an integer and
$g^2-\Delta c^2=1$.  Writing
\[
 \chi_A(r)+\psi_A(r)\sqrt{A^2-1}
       =(A+\sqrt{A^2-1})^r,
\]
Pell classification gives $c=\psi_{a+2}(p)$ and
$k=2\psi_{2XY^2+1}(m_1)$.  Here $p\ge2$ and $m_1\ge2$:
$m_1=1$ would give $k=2$, $c=2Y+1<2(a+2)$.
Consequently, still before proving $C>0$,
\begin{equation}
 c>kY\ge4(2XY^2+1)Y>8XY^3\ge8q^9.
 \label{eq:budget-c-growth}
\end{equation}

\subsection{The new conditional growth bound}

\begin{lemma}[Strong-Pell growth budget]
Suppose $q\ge2$, $Y\ge q^3$, $A\ge2Y+2$,
$\Delta=A^2-1$, $c=\psi_A(p)\ge6$, and positive integers $f,i$
satisfy $f^2-\Delta(ic^2)^2=1$.  Then
\[
 i>8q\Delta c^4,\qquad i^3>6qf^2,\qquad i^3>q^2,
 \qquad Aic^2<i^2.
\]
If $0<k<c$, also $i^3>Aic^2+k$.
\end{lemma}
\begin{proof}
Classify $f=\chi_A(m)$ and $ic^2=\psi_A(m)$.
Strong divisibility gives $p\mid m$: the addition law and
$\gcd(\chi_A(r),\psi_A(r))=1$ give the Euclidean identity
$\gcd(\psi_A(r),\psi_A(s))=\psi_A(\gcd(r,s))$.
Write $m=pv$ and $g=\chi_A(p)$.  Composition gives
$ic=\psi_g(v)$.  As $g^2\equiv1\pmod c$, recurrence gives
$\psi_g(v)\equiv vg^{v-1}\pmod c$; hence $c\mid v$.
Thus $v\ge c\ge6$ and $ic>g^5$.  Since
$g^2=1+\Delta c^2$ and $\Delta>4Y^2$,
\[
 i>\Delta^{5/2}c^4
   >8Y^3\Delta c^4\ge8q^9\Delta c^4\ge8q\Delta c^4.
\]
Putting $H_0=\Delta c^4\ge15$ gives
$i^3>8qH_0i^2>6q(H_0i^2+1)=6qf^2$.
Also $i>q$, $i>Ac^2$, and $i>c>k$.
Therefore $Aic^2<i^2$ and $i^3-i^2>k$.
\end{proof}

Apply the lemma after \eqref{eq:budget-radical-recovery}, with
$A=a+2$ and $f^2=D_f$.  It gives
\begin{equation}
 i^3>6qD_f,\qquad i^3>q^2,\qquad i^3>\bar f+k.
 \label{eq:budget-bounds}
\end{equation}
Also $\bar f>f$, since $\bar f^2-D_f=I^2-1>0$.
These are consequences of the recovered Pell equations; no assertion
of these bounds on arbitrary positive inputs is needed.

\subsection{The noncircular positivity bootstrap}

Suppose $C\le0$, put $d=-C$ and $B=-B_t>0$.  Then
\[
 B=(qK_0+X)d+qF,\qquad
 R=[q(q-F)+d](q^2-1)+(M_C+qM_F)J.
\]
The mask bounds imply
\begin{equation}
 R<qB+q^4+q^2,\qquad R<0\Longrightarrow -R<q^2B.
 \label{eq:budget-rank-bounds}
\end{equation}
Let $s=Y/q^3>0$.  Equation \eqref{eq:budget-radical-recovery} becomes
\[
 \tau_*+f_*+y_*
  =\ell+k+n+cT+V+\bar f+s+\frac{WB}{q-1}.
\]
Its left side is at most $t+f+y$, with $t<\ell+k$.
We derive a contradiction in both possible signs of $V$.

As $\delta>D_f$,
\[
 0<y-|V|=\frac{V^2-1}{\delta(y+|V|)}<\frac{|V|}{\delta}.
\]
If $V>0$, the integer $b=cT-R$ is positive, so $y-V<b$ and
\[
 cT+V+\bar f-f-y>cT-b=R.
\]
For $R\ge0$ the right side of the radical equality already exceeds
$t+f+y$.  For $R<0$, the remaining term exceeds
$h_*B\ge i^3B>q^2B>-R$, with the same contradiction.

If $V<0$, then $R\ge cT\ge c>8q^9$.  Thus
$q^4+q^2<R/2$ and \eqref{eq:budget-rank-bounds} gives $R<2qB$.
We have $|V|<D_fR$ and $y-|V|<R$, whence
\[
 cT+V+\bar f-f-y>-(2D_f+1)R>-3D_fR.
\]
But $WB/(q-1)>h_*B\ge i^3B>6qD_fB>3D_fR$, again a
contradiction.  The case $V=0$ was excluded earlier.  Therefore $C>0$.

Now $3q<R<q^4-2$, and \eqref{eq:budget-c-growth} gives $c>2R$.
Hence $V>1$ and the positive-gap estimate strengthens to
\[
 cT+V+\bar f-f-y>R>0.
\]
Restore positive source witnesses by
\[
 f_{\rm old}=f,\qquad u=y-V,\qquad
 r=n+cT+V+\bar f-f_*-y_*.
\]
Here $u>0$ and $r>n>0$ for every permitted choice of signs.
Subtracting $M(f_*+y_*)$ from $\bar Z$ gives exactly the first source
carrier $Z=M\tau_*$.  All four source equations hold, so source
soundness yields a finite ZFC refutation.

Conversely, choose canonical outer, strong, and auxiliary source
witnesses first, then set $h_*=i^3+cT$.
By \eqref{eq:budget-bounds}, $h_*>cT+\bar f+k$; the weighted
source construction supplies $r>qW-2k>cT+\bar f$.
Take
\[
 n=r-cT-\bar f+f+u>0.
\]
Then $\bar Z=M(\tau+f+y)$, so $N_8=0$, and the unchanged canonical
main carrier gives $A_1=0$.  No new congruence on the strong pair is
required.  The witness choices have no dependence loop.
Composing both implications with the inherited source calibration proves
\[
 \mathrm{ZFC}\vdash
 \bigl(\nexists\boldsymbol{x}\in\Zp^{11}\;
       P_{11}(\boldsymbol{x})=0\bigr)\ \longleftrightarrow\ \ConZFC.
\]

\subsection{Exact degree and reproducible finite evidence}

Put $e=I^2-1$ and $z=\bar Z-M(V+\bar f)$, so
$D_f=\bar f^2-e$ and $\deg z=9$.
Expand the universal norm using
\[
 x=M(V+t)+z,\quad A=M^2D,\quad B=M^2(t^2-e),\quad
 C_{\rm num}=M^2(QV^2-1),\quad\delta=Q-1.
\]
With weights $(18,18,4,7,9,10,8)$ on $(Q,V,M,t,z,e,D)$,
the exact $1110$-term integer expansion has degree $200$ and leading form
\begin{equation}
 64M^6Q^4V^4t^2z^2.
 \label{eq:budget-leading-norm}
\end{equation}
The checker expands this identity exactly; naive gate-degree propagation
does not suffice.  Write $q_0=3\beta j$, $k_0=\eta+\zeta$, and
\[
 B_{t,0}=(K_0q_0+X)(q_0-F-\alpha)-q_0F.
\]
The relevant leading forms are
\begin{align*}
 D_{f,0}&=i^2k_0^4X^2Y^6,&Q_0&=i^2k_0^4X^4Y^8,\\
 R_0&=q_0^3(q_0-F),&V_0&=-D_{f,0}R_0,\\
 \bar f_0&=ik_0^2XY^3,&M_0&=q_0^4,\\
 z_0&=-q_0^4(i^3+k_0YT)B_{t,0},\\
 (A_1)_{13}&=32q_0^2k_0X^2Y^4
       \{4q_0^2Y\gamma+q_0^3(q_0-F)\}.
\end{align*}
Therefore $\deg P_{11}\le213$.  Its coefficient at
\[
 j^{48}i^{24}\eta^{37}X^{30}Y^{67}F^4\alpha^2\gamma
\]
is exactly $2^{13}(3\beta)^{48}>0$.  The minimum $q_0$ degree and
the maximum displayed $i$ degree select this coefficient uniquely in
\eqref{eq:budget-leading-norm} times $(A_1)_{13}$; it is independent of
$K_0,M_C,M_F$.  Thus the exact degree on the actual numeral tuple is
$213$.

The standard-library checker passes normally and with Python assertions
disabled.  It checks the complete prefix, the octic conjugate identity,
the weighted leading form, the nonzero degree certificate, coefficient
hashes, deterministic rebuild, all eleven live inputs, and all $112$
live arithmetic operations.  Constructed Pell examples supplement the
unbounded proof above; they do not test consistency or replace the
inherited semantic calibration.

\section{Seventeen integers: compressing congruence witnesses}
\label{sec:integer17}

The following constructions use unrestricted integer coordinates. They give
the variable--degree pairs $(23,10)$, $(17,88)$, and $(16,184)$,
with a $(16,192)$ sum-of-squares control.
The first two improve the incoming $(24,10)$ and $(19,88)$ points
\cite{budget,joint}; the last trades degree for another coordinate.
Their literal circuits use respectively $131$, $114$, and $135$ binary
operations; the degree-192 control uses $134$. Scalar arithmetic and
squaring are counted. These are different
tradeoffs from the much higher-degree ten-integer reductions. In particular,
the $114$-operation circuit is larger than the incoming $109$-operation
nineteen-integer circuit. No global minimum or historical priority is claimed.

All four use the same six fixed integers
\[
 \beta=2^{1064921}-1,\qquad K_0,M_C,M_F,\qquad
 A_c=2^{1113},\quad B_c=2^{666104}.
\]
The first four are the entire source-34
\texttt{mixed-parity-unrounded-1064921} tuple; $M_F$ already includes its
shift by $\beta$. The last two are source71's determinant scales
\cite{source71}. The actual six integers, compressed minimal little-endian
bytes, finite native recipe, and complete raw hashes are supplied in the
integer component's \texttt{coefficients/} directory. The raw encodings
total $615\,883$ bytes. These fixed leaves are not additional unknowns.
They satisfy
\[
 \beta\ge17,\quad\beta\equiv1\pmod3,\quad
 0<M_C<\beta<M_F<2\beta,\quad K_0\ge0,\quad A_c,B_c\ge2.
\]

\subsection{Two inherited guards}

Put $\operatorname{Pos}(p_1,p_2,p_3)=1+p_1^2+p_2^2+p_3(p_3+1)$.
It takes exactly the positive integer values, by the three-square theorem
applied to $4n-3$. The nearest-square joint-positivity lemma used in the
incoming constructions is
\begin{equation}\label{eq:integer-pair}
 x,y>0\quad\Longleftrightarrow\quad
 \exists t,p_1,p_2,p_3\in\Z:\quad
 xy-(t^2-x-y+1)^2=\operatorname{Pos}(p_1,p_2,p_3).
\end{equation}
Indeed a positive left side forces $xy>0$ and excludes both inputs being
negative. Conversely, for $x,y>0$, choose the nearer neighboring square
to $x+y-1$. Its squared error is less than $x+y-1\le xy$, leaving a
positive integer represented by $\operatorname{Pos}$.

For $\alpha=q-F-C$ define
\[
 d=(A_cF)(B_cC)-(q-A_cF-B_cC)^2.
\]
The inequality $d>0$ forces $F,C,\alpha,q$ to have one common nonzero
sign. To see this, put $U=A_cF$, $B=B_cC$. Then $U,B$ have a common
sign and $|q-U-B|<\sqrt{UB}\le(|U|+|B|)/2$. Hence $q$ has their sign
and $|q|>(|U|+|B|)/2\ge|F|+|C|$, which also fixes the sign of
$\alpha$. Source71's separate canonical-witness estimates prove
completeness for this determinant; that calibrated existence restriction
is inherited, not re-proved by finite sign tests.

\subsection{The complete polynomial}

The seventeen independent integer coordinates are
\[
 (J,C,z,f,h,r,v,w,\tau,\nu,\gamma,T,c,t,p_1,p_2,p_3).
\]
The following are polynomial abbreviations only:
\begin{align*}
 J_0&=3J+1,&q&=\beta J_0+1,&W&=4(3w+1),\\
 X&=qW,&Y&=(2v+1)q^3,&E&=XY,\\
 a&=(X+1)Y,&H&=4a+3,&\Delta&=a^2+H,\\
 F&=(K_0+W)C-(q-1)z,&\alpha&=q-F-C.&&
\end{align*}
Continue with
\begin{align*}
 R&=(q-F)(q^3-q)-C(q^2-1)+(M_C+qM_F)J_0,\\
 k&=R+1+hE,\qquad \eta=c-Yk,\qquad \zeta=k-\eta,\\
 I&=(2r+1)c^2,\qquad V=c(Tf-1)-Rf^2,\\
 m&=XY^2,\qquad B_0=mk,\qquad s_0=X+\gamma H,\qquad Z=\eta\zeta.
\end{align*}
Define five residuals:
\begin{align}
 A_1&=\tau^2+(2\tau-k)B_0-1,\nonumber\\
 A_2&=2acs_0+s_0^2-Hc^2-1,\nonumber\\
 A_3&=I^2-\Delta(f^2-1),\label{eq:integer17-rows}\\
 A_4&=(V+\nu)^2-I^2\nu(2V+\nu)-1,\nonumber\\
 A_5&=dZ-(t^2-d-Z+1)^2-\operatorname{Pos}(p_1,p_2,p_3).\nonumber
\end{align}
The ordinary integer polynomial is
\begin{equation}\label{eq:integer17}
 \boxed{P_{17}=A_1^2+A_2^2+A_3^2+A_4^2+A_5^2.}
\end{equation}

The centered roots $\tau_{\rm old}=\tau+B_0$ and
$y_{\rm old}=\nu+V$ give the identities
\[
 A_1=\tau_{\rm old}^2-m(m+1)k^2-1,\qquad
 A_4=I^2(V^2-y_{\rm old}^2)+y_{\rm old}^2-1.
\]
Also $D=ac+s_0$ gives $A_2=D^2-\Delta c^2-1$.

\subsection{One quotient replaces two}

\begin{lemma}[Integer quotient compression]\label{lem:integer-quotient}
Suppose $c\ne0$ and $c\mid f^2-1$. The simultaneous integer identities
\[
 V=of-c=jc-R
\]
are equivalent, after projection of $j,o$, to
$V=c(Tf-1)-Rf^2$ for one integer $T$.
\end{lemma}
\begin{proof}
From the old identities, $c\mid of+R$. The identity
\[
 f(of+R)-o(f^2-1)=o+Rf
\]
shows that $T=(o+Rf)/c$ is integral and has the desired value of $V$.
Conversely put $b=(f^2-1)/c$ and define
\[
 o=cT-Rf,\qquad j=Tf-1-Rb.
\]
These are integers and restore both identities. The witness-restoration
division introduces neither division nor a $b$ coordinate into $P_{17}$.
\end{proof}

The order in which $c\mid f^2-1$ is established matters. Invoking an
auxiliary-index theorem that already used the deleted congruences would
be circular. Here the needed bootstrap is source70's signed orientation
argument \emph{before} its auxiliary-congruence step
\cite{source70}. We spell out its dependency boundary.

At a $P_{17}$ zero, $A_5=0$ gives $d>0$ and $\eta\zeta>0$.
Write $Q=|q|$, $U=|X|$, $Z_0=|Y|$, $K=|k|$, and $M=|E|$.
The determinant, authentic mask bounds, and affine residue restrictions give
\[
 Q\ge16,\quad Q(Q^2-1)<R<Q^4-2,\quad
 U\ge4Q,\quad Z_0\ge Q^3,\quad M>3R+1.
\]
The positive product gives $K\ge2$, $0<|\eta|<K$, and
$K(Z_0-1)<|c|<K(Z_0+1)$. Let $B=|a+2|$ and $P=|2XY^2+1|$.
They satisfy $P>B\ge2$ and $B>R$. Classifying the first and main
Pell norms gives
\[
 K=2\psi_P(n_0),\qquad |c|=\psi_B(p),\qquad p>n_0.
\]
The recurrence modulo $M$ and $k=R+1+hE$ give
$\pm2n_0\equiv R+1\pmod M$. The positive branch gives
$n_0\ge(R+1)/2$, while the negative branch gives $n_0>R$.
Thus $p>n_0\ge24$; Pell growth yields
\[
 |c|>B(B^2-1)^2,\qquad |c|>2p,\qquad |c|>2R.
\]
The inherited ordinary strong-rank lemma can now be applied to $A_3$.
It gives $|f|=\chi_B(l)$ with $p\mid l$, without using either missing
congruence. Therefore $\psi_B(p)\mid\psi_B(l)$ and
\[
 c\mid (B^2-1)\psi_B(l)^2=f^2-1.
\]
The precise retained source proof is
\texttt{upstream/ORIENTATION\_PROOF.md}, Sections 3--4.2; its
auxiliary-congruence argument starts in Section 4.3.

Lemma~\ref{lem:integer-quotient} restores $j,o$. All other deleted
coordinates are monic polynomial definitions, the two root shifts are
integral, and the old positive guard triples can be restored by
$\operatorname{Pos}$. This gives a zero of the incoming calibrated
37-integer affine quartic \cite{round11}. Conversely, any zero of that
source supplies the same divisibility, hence $T$; its transport row gives
$F=(K_0+W)C-(q-1)z$ after the invertible change
$\alpha=q-F-C$. Forget the monic coordinates, perform the root shifts,
and use \eqref{eq:integer-pair}. Thus $P_{17}$ has an integer zero
exactly when that source does. This is a new existence equivalence on
its full integer solution set, with no additional canonical restriction.

The exact verifier restores all 22 arithmetic rows of the literal parent:
four match the new arithmetic residuals, one is $\pm R(cb-f^2+1)$,
and all others vanish identically. It thereby exposes the sole use of the
proved divisibility rather than silently treating $b$ as an unknown.

\section{Sixteen integers: merging an ordinary and a rational Pell root}
\label{sec:integer16}

\begin{lemma}[Separated rational norm]\label{lem:integer-rational-norm}
Let $A,Q,S,V\in\Z$, with $A>0$, $Q>1$, and $S^2>A$.
Put $\delta=Q-1$ and $B_{\rm num}=S^2(QV^2-1)$. Then
\[
 \exists n\in\Z:\quad
 [\delta(n^2-A)-B_{\rm num}]^2-4\delta A B_{\rm num}=0
\]
if and only if there are integers $u,y$ with
$u^2=A$ and $QV^2-(Q-1)y^2=1$.
\end{lemma}
\begin{proof}
If $V=0$, the displayed expression is a square plus the strictly positive
integer $4\delta AS^2$. Thus any zero has $|V|\ge1$. The rational number
$B=B_{\rm num}/\delta$ then satisfies $B\ge S^2>A$.
Dividing by $\delta^2$ gives $(n^2-A-B)^2=4AB$.
Here $n\ne0$, since $n=0$ would give $A=B$. The identity
\[
 (n^2+A-B)^2-4n^2A=(n^2-A-B)^2-4AB
\]
shows that $u=(n^2+A-B)/(2n)$ is a rational square root of the integer
$A$, hence an integer. The rational $s=n-u$ satisfies $s^2=B$, so
$y=s/S$ is rational and obeys $QV^2-(Q-1)y^2=1$.

Source69's adjacent-coefficient Pell descent makes this $y$ integral.
For completeness, if $y=a/d$ is reduced with $d>0$, the norm gives
$d^2\mid Q-1$. Make $V,y$ positive by changing signs. Unless $V=1$,
the integral coordinate
\[
 V'=(2Q-1)V-2(Q-1)y
\]
and rational $y'=(2Q-1)y-2QV$ preserve the norm. The reduced denominator
of $y'$ remains $d$, since $2Q-1\equiv1\pmod d$.
The inequalities $y>V$ and $QV>(Q-1)y$ imply $-V<V'<V$; the norm
excludes $V'=0$. Strong induction on $|V|$ reaches $V=1$, where
$y=1$ and $d=1$. This is the inherited descent proof, applied to the
new separated norm. Conversely, take $n=u+Sy$ and factor the norm.
\end{proof}

Delete $\tau,\nu$ from the seventeen coordinates and add one unrestricted
integer $n$. Keep every other abbreviation above and put
\begin{align*}
 A_0&=1+m(m+1)k^2,\qquad S=k(2m+1),\qquad\delta=I^2-1,\\
 L&=\delta\bigl[n^2+2n(B_0+SV)+2B_0SV-mk^2-1\bigr]
       +S^2(1-V^2),\\
 N&=L^2-4\delta A_0S^2(I^2V^2-1).
\end{align*}
The new equation is the four-square polynomial
\begin{equation}\label{eq:integer16}
 \boxed{P_{16}=A_2^2+A_3^2+A_5^2+N^2.}
\end{equation}
It is exactly the norm of Lemma~\ref{lem:integer-rational-norm} with
$Q=I^2$ and $n_{\rm total}=n+B_0+SV$, because
\[
 \delta(n_{\rm total}^2-A_0)-S^2(I^2V^2-1)=L.
\]
This centering cancels the two leading $V^2$ terms and lowers the degree.

\paragraph{The separation does not assume either missing norm.}
From $A_5=0$ alone, $d>0$ and $\eta\zeta>0$. The determinant and
fixed coefficient bounds imply $|q|\ge16$; the definitions give
$|Y|\ge|q|^3$ and $|X|\ge4|q|$. As $\eta,\zeta$ have the same
sign, $|k|\ge2$ and
\[
 |c|>|k|(|Y|-1)>1.
\]
Hence $I^2>1$, while $|m|>1$ gives $A_0>0$. Finally the exact identity
\begin{equation}\label{eq:integer-separation}
 S^2-4A_0=k^2-4\ge0
\end{equation}
proves $S^2>A_0$. All hypotheses of the rational-norm lemma are therefore
established before either deleted Pell root is extracted.

At a zero of $P_{16}$, the lemma supplies integral $\tau_{\rm old}$
and $y_{\rm old}$; put $\tau=\tau_{\rm old}-B_0$ and
$\nu=y_{\rm old}-V$. This restores a zero of $P_{17}$. Conversely a
$P_{17}$ zero gives a $P_{16}$ zero by $n=\tau+S\nu$.
This proof covers arbitrary signed coordinates, and introduces no new
divisibility restriction on canonical witnesses.

\subsection{Exact degree certificates}

For $P_{17}$ the degrees of $X,Y,E,a,R,k$ are $2,4,6,6,5,7$;
$\eta,\zeta$ have degree 11 and $d$ has degree 4.
The five residual bounds are $24,14,14,14,44$.
Only the joint guard reaches degree 44. Its top part is
$-\eta_{\rm top}^4$, with
\[
 \eta_{\rm top}=-2^4 3^8\beta^7hJ^7wv^2.
\]
Consequently
\[
 [h^8J^{56}w^8v^{16}]P_{17}=2^{32}3^{64}\beta^{56}>0,
\]
proving exact degree 88 on the actual fixed coefficient slice.

For $P_{16}$, $\deg S=17$ and $\deg V=7$.
The top part of $L$ is $-S_{\rm top}^2V_{\rm top}^2$, of degree 48;
the second term of $N$ has degree at most 94. Thus $N$ has degree 96.
Specifically,
\begin{align*}
 S_{\rm top}&=2^8 3^{13}\beta^{11}hJ^{11}w^2v^3,\\
 V_{\rm top}&=f^2(324\beta^3J^3wC-81\beta^4J^4z),
\end{align*}
and
\[
 [h^8J^{120}w^{16}v^{24}z^8f^{16}]P_{16}
   =2^{64}3^{136}\beta^{120}>0.
\]
The other squared rows have degrees at most $28,28,88$, proving exact
degree 192. The direct uncentered weighted norm would have upper degree
216. Exact homogeneous-component arithmetic, rather than a naive gate
bound, verifies both coefficients in the supplied auditor.

\subsection{A degree-184 finalizer}

Let $G=A_2^2+A_3^2+A_5^2$. Instead of $G+N^2$, define
\begin{equation}\label{eq:integer16unit}
 \boxed{P_{16}^{\rm unit}=G+N+GN=(1+G)(1+N)-1.}
\end{equation}
For every integer assignment $G\ge0$. If the new polynomial is zero,
the positive integer $1+G$ multiplies the integer $1+N$ to give 1;
therefore $G=N=0$. Conversely these equalities make both polynomials
zero. This exact zero-set equivalence needs no sign or congruence
assumption about $N$.

The degree is $88+96=184$, certified by
\[
 [h^{12}J^{116}w^{16}v^{28}z^4f^8]P_{16}^{\rm unit}
   =2^{64}3^{132}\beta^{116}>0.
\]
The expanded form $G+N+GN$ uses $135$ binary operations, one more than
the degree-192 control. The archive retains both explicit polynomials.

\section{A lower-degree companion in twenty-three integers}
\label{sec:integer23}

This companion combines the incoming residue orientation with the paired
Pell-root translations \cite{round11,budget}. Its independent coordinates are
\[
 \begin{gathered}
 J,F,\alpha,z,f,h,r,v,w,\tau,\eta,\zeta,\nu,\gamma,o,j,\ E,a,I,\ t,p_1,p_2,p_3.
 \end{gathered}
\]
Here $F,\alpha,E,a,I$ are coordinates, not their expressions in the
seventeen-variable equation. Define
\begin{align*}
 J_0&=3J+1,&q&=\beta J_0+1,&W&=4(3w+1),\\
 X&=qW,&Y&=a-E,&k&=\eta+\zeta,\\
 c&=Yk+\eta,&V&=fo-c,&C&=q-F-\alpha,\\
 H&=4a+3,&\Delta&=a^2+H,&B_0&=E(c-\eta),\\
 s_0&=X+\gamma H.&&&&
\end{align*}
Use
\[
 R=((q-1)(q-F)+\alpha)(q^2-1)+(M_C+qM_F)J_0
\]
and the earlier definitions of $d$, $Z=\eta\zeta$, and $A_5$.
The eleven residuals, whose squared sum defines $P_{23}$, are
\begin{align*}
 &a-E-(2v+1)q^3, && E-X(a-E),\\
 &\tau^2+(2\tau-k)B_0-1,&&2acs_0+s_0^2-Hc^2-1,\\
 &k-1-hE-R,&&(K_0+W)C-(q-1)z-F,\\
 &V+k-1-cj-hE,&&I-(2r+1)c^2,\\
 &I^2-\Delta(f^2-1),&&(V+\nu)^2-I^2\nu(2V+\nu)-1,\\
 &A_5.&&
\end{align*}
The count is $37-6+4-12=23$: replace the source affine quartic's six
private guard coordinates by four, translate both roots, and delete twelve
monic coordinates. Their exact restorations are
\begin{align*}
 N&=q^3,&X&=qW,&L&=q^2-1,&C_2&=c^2,\\
 F_2&=f^2,&B_{\rm aux}&=I^2,&g&=\nu(2V+\nu),&\Delta&=a^2+4a+3,\\
 D&=ac+s_0,&P&=k+B_0,&c&=(a-E)k+\eta,&V&=fo-c,
\end{align*}
with $\tau_{\rm old}=\tau+B_0$, $y_{\rm old}=V+\nu$.
Every arithmetic map is integral and polynomial. The auditor checks all
22 source arithmetic rows: twelve vanish identically and ten match the
displayed arithmetic rows once each, up to sign. It also authenticates
both old guard arguments.

The residual degrees are at most
$4,3,4,5,4,2,3,5,4,5,4$. Only the $I$ definition contains $r$, so
\[
 [r^2a^4\eta^4]P_{23}=4.
\]
Thus the degree is exactly ten, independently of the fixed ports.

\begin{theorem}[Inherited calibration, new integer reductions]
For each $P\in\{P_{23},P_{17},P_{16},P_{16}^{\rm unit}\}$, composition of the new
reductions above with the named source native/compiler, Pell, and
syntactic-refutation calibration theorems proves
\[
 \mathsf{ZFC}\ \vdash\
  \bigl(\neg\exists\mathbf x\in\Z^{\dim P}:P(\mathbf x)=0\bigr)
  \ \longleftrightarrow\ \ConZFC.
\]
\end{theorem}

The theorem asserts the equivalence, not either side. The full native and
logical calibration is inherited; the finite algebra/data checks are not an
end-to-end Lean/Rocq or encoded internal-ZFC certificate. All four literal
packets, inlined expressions, actual coefficients, exact degree receipts,
source restorations, and shared DAGs are included. The verifier uses only
Python's standard library and works with assertions disabled. This component's
new steps have independent mathematical review within the AI collaboration;
that is not external human peer review.


\section{Reproducibility, audit coverage, and repository use}
\label{sec:reproduction}

\subsection{Complete literal artifacts}

The JSON packets encode finite division-free arithmetic DAGs with a
single output. The integer component additionally includes fully
inlined residual expressions. The variable lists are part of the packet;
temporary registers do not enlarge them. All large fixed coefficients
are present as compressed minimal little-endian integer bytes, with
their signs where needed and SHA-256 bindings. The nine-positive
majorant is a supplied integer coefficient table rather than a runtime
Laurent-series operation.

The principal files are:
\begin{center}
\small
\begin{tabular}{p{64mm}p{76mm}}
\toprule
Purpose & Path\\
\midrule
Nine-positive literal polynomials & \path{nine/equations/}\\
Majorant and positive defect tables & \path{nine/data/}\\
Independent nine-positive proof and audit & \path{nine/certificates/}\\
Eleven-positive polynomial & \path{positive11/p11_positive_degree213.json}\\
Integer literal circuits & \path{integer/circuits/}\\
Integer inlined residual packets & \path{integer/equations/}\\
Expanded ordinary proofs & \path{positive11/PROOF.md}, \path{integer/PROOF.md}\\
Input provenance and formal status & \path{provenance/}, \path{PROOF_STATUS.md}\\
Release integrity and fresh transcripts & \path{MANIFEST.sha256}, \path{verification/}\\
\bottomrule
\end{tabular}
\end{center}

\subsection{Exact checks and diagnostics}

The release driver \file{verify_all.py} invokes five audits in normal
Python and with assertions disabled by \file{-O}. All active checks use
explicit error conditions rather than relying on removable assertions.
The audits have the following roles.

\begin{itemize}
 \item The primary nine-positive audit verifies the literal regenerated
 circuits, source extraction, actual fixed coefficient recipes,
 topological order, liveness, operation counts, and cancellation-aware
 degrees. Independently written direct modular formula evaluations are
 transcription diagnostics.
 \item The independent H1868 audit imports no producer. It reconstructs
 all 22 source identities and 20,466 convolution interactions, the
 source height and binary population threshold, the 3035-term radicand,
 the 333-term corrected majorant, the 36,309-term positive shifted
 defect, and the unique highest-degree contribution.
 \item The local saturation audit checks 5472 arithmetic cases and
 constructs three complete Pell fixtures. These are diagnostic
 controls for the separate unbounded proof.
 \item The eleven-positive audit binds the exact circuit to the literal
 formulas, checks fixed coefficient bytes, expands the 1110-term
 double-centered norm in abstract variables, and verifies a nonzero
 degree-213 coefficient. Constructed Pell fixtures test the growth
 mechanism without standing in for its proof.
 \item The integer audit checks all source arithmetic restorations,
 the sole quotient-divisibility remainder, exact local norm identities,
 rational-root diagnostic cases, degree-attaining coefficients, fixed
 ports, and the complete literal DAGs. It does not assume degree
 attainment from an upper-bound propagation.
\end{itemize}

The positive majorant certificate is stronger than a collection of sampled
inequalities: its exact translated polynomial identity and positive
coefficients prove the bound throughout the stated domain. Conversely,
the modular evaluations and small Pell examples are not represented as
proofs of the compiler semantics or of ZFC consistency.

From the root of an extracted archive, the full finite replay is
\begin{verbatim}
python3 verify_all.py --both
\end{verbatim}
Only Python 3.10 or later and its standard library are needed. No symbolic
algebra package is required for this replay. The generator and each
verifier can also be run separately; the root README gives the exact
commands. Compilation of this report needs only a conventional TeX Live
installation:
\begin{verbatim}
latexmk -pdf -interaction=nonstopmode -halt-on-error report.tex
\end{verbatim}

\subsection{Integration and limitations}

The archive is suitable for placement as another research contribution
under the repository's Diophantine-consistency collection. It need not
replace any currently certified equation or change an existing theorem
statement. A reviewer can inspect the ordinary lemmas, replay the finite
checks, and select an equation according to the domain and cost measure
that matters for the next stage.

The independently implemented audits and reciprocal mathematical reviews
were performed within the AI collaboration. They are not external human
peer review. The finite data are authenticated against the supplied
snapshot; the complete older native/compiler and logical-calibration
chain remains an identified dependency. The number of variables, exact
degrees, circuit counts, and coefficient payloads are concrete artifacts.
Their global optimality is not established.

\section{Alternative encodings and what remains open}
\label{sec:alternatives}

The request permits constructions unrelated to K215 or a Jones universal
polynomial. We examined that possibility as a separate design question.
The repository already has a direct annotated checker for sealed ZFC
proofs, explicit primitive-recursive programs for its syntax operations,
and a complete numeric refutation checker \cite{repo}. Merely suggesting
direct proof encoding would therefore repeat existing work.

A direct checker could provide a valuable route to the exact internal
ZFC target without importing the numerical machine-calibration bridge.
However, a checker program, an existence theorem from MRDP, and a
fully instantiated small ordinary polynomial are different deliverables.
To compete with the present equations, that route still needs a measured
shared arithmetic graph, a complete explicit coefficient binding, and
unbounded source/target correctness at the chosen witness domain.
An unspecified program-index constant does not resolve those obligations.

The active repository reports a complete direct-checker program tree of
83,281 nodes, or 76,313 after its certified optimization, while its
smaller shared quadratic graph remains unfinished \cite{repo}.
Those are program-tree sizes, not counts of Diophantine variables.
We have not converted them into an advertised smaller equation.
Likewise, the ten-integer results of Sun are relevant comparison
architecture \cite{sun}, but do not by themselves materialize a smaller
fixed consistency polynomial for this particular calculus.

The useful general lesson from the new constructions is local and concrete:
when an arithmetic certificate already forces a square or a divisibility,
its polynomial presentation can sometimes exploit that consequence
without preserving the original witness parameterization. The soundness
proof must establish the consequence before using it; completeness can
rebuild suitable witnesses existentially. Positive Pell saturation,
quotient recovery, and the guard-only root separation each apply that
principle at a different stage. They improve completed equations while
leaving the compiler-independent route as an explicit open construction
problem.

\begin{thebibliography}{99}
\bibitem{repo}
Vladimir Reshetnikov and ProveIt contributors.
\emph{ProveIt: DiophantineConsistency research and incoming contributions}.
Reviewed commit
\texttt{c1a102936aca437e9999\allowbreak530a0111e6ddab56dc6d},
24 October 2026.
\href{https://github.com/VladimirReshetnikov/ProveIt/tree/c1a102936aca437e9999530a0111e6ddab56dc6d/Computability/DiophantineConsistency}{Pinned research directory}.
The source status, size plan, K215 status, and object-ZFC obligation are
retained under \file{provenance/}.

\bibitem{round11}
ProveIt incoming contribution.
\emph{Structural Reductions, Round 11}.
\file{ConZFC_Structural_Reductions_Round11_2026-10-24.zip},
at the pinned repository's \file{docs/incoming/}.
H1868 packing, norm saturation for the ten-integer construction, and
the 37-integer affine quartic source.

\bibitem{joint}
ProveIt incoming contribution.
\emph{ZFC Diophantine Joint Guards}.
\file{ZFC_Diophantine_Joint_Guards_2026-10-24.zip}.
Joint positivity, integer low-degree alternatives, and the
eleven-positive centered interfaces.

\bibitem{budget}
ProveIt incoming contribution.
\emph{Budget and Centering}.
\file{con-zfc-budget-and-centering-2026-10-23.zip}.
Paired root translations and integer degree-ten comparison.

\bibitem{radical}
ProveIt incoming contribution.
\emph{Radical Guard Descent / Cross-Guard Descent}.
\file{con-zfc-radical-guard-descent-2026-10-23.zip}.
Whole-radical majorants and coefficient/data recipes.

\bibitem{reportincoming}
ProveIt incoming contribution.
\emph{Standalone report}, \file{report.zip}, containing \file{report.tex},
at the pinned \file{docs/incoming/}.
Its nine-positive degree-\(1\,158\,240\), 811-operation entry is a reported
comparison; this archive contains no executable companion.

\bibitem{source54}
ProveIt contributors.
\emph{54-calibrated-carry} and the later compact54-sharp H2082 candidate.
The complete coherent compact tuple, source adapter, inherited proof
notes, and 440-operation parent DAG are retained under
\file{nine/baseline/compact2082/}.

\bibitem{source70}
ProveIt contributors.
\emph{70-residue-orientation}.
Signed strong-rank recovery and the eleven-positive native Pell source.
Relevant proof files are retained under \file{integer/upstream/}
and \file{positive11/provenance/}.

\bibitem{source71}
ProveIt contributors.
\emph{71-determinant-guards}.
Positive nine-coordinate and calibrated determinant interfaces.
Relevant proofs, scales, and notices are retained in the corresponding
components. Source labels follow the pinned repository placements.

\bibitem{productpell}
ProveIt contributors.
\emph{ProductPell.lean}, the positive Pell certificate in the formalized
Jones development.
\href{https://github.com/VladimirReshetnikov/ProveIt/blob/c1a102936aca437e9999530a0111e6ddab56dc6d/Computability/HilbertTenthProblem/Lean/Diophantine/Paper1982/ProductPell.lean}{Pinned source file}.

\bibitem{jones}
James P. Jones.
Universal Diophantine equation.
\emph{The Journal of Symbolic Logic} \textbf{47}(3), 549--571, 1982.

\bibitem{sun}
Zhi-Wei Sun.
\emph{Ten unknowns for Hilbert's tenth problem over the integers}.
arXiv:2609.01594, version 2, 2 September 2026.
\url{https://arxiv.org/abs/2609.01594}.

\bibitem{parallel}
Jonas Bayer, Marco David, Malte Ha{\ss}ler, Yuri Matiyasevich,
and Dierk Schleicher.
\emph{Diophantine equations over \(\mathbb Z\): universal bounds and
parallel formalization}. arXiv:2506.20909, version 2, 2025.
\url{https://arxiv.org/abs/2506.20909}.
\end{thebibliography}

\end{document}
